\section{Tail Inference}
\label{app:inference}

\paragraph{Same-specification standard-error ratio.}
The case for bootstrap inference at the tail
(Section~\ref{sec:method-inference}) rests on a same-specification
comparison. We take the bootstrap standard error from the exact headline-spec
bootstrap and the kernel standard error from the same table, at
$\tau = 0.01$ (Table~\ref{tab:seratio}). The ratio ranges from
\SeRatioMin{} to \SeRatioMax{} across horizons. This compares like with like.
A cross-specification comparison would instead conflate the standard-error
gap with differences in scale, controls, and shock definition.

\begin{table}[t]
\centering
\caption{Kernel versus bootstrap standard errors, same specification,
$\tau = 0.01$.}
\label{tab:seratio}
% GENERATED by scripts/paper/make_tables.py — DO NOT EDIT
% fragment only: wrap in \begin{table} + caption in the manuscript
\begin{tabular}{lccc}
\toprule
$h$ & kernel SE & bootstrap SE & ratio \\
\midrule
0 & 0.0057 & 0.0094 & 1.65 \\
3 & 0.0229 & 0.0362 & 1.58 \\
6 & 0.0344 & 0.0931 & 2.71 \\
12 & 0.0404 & 0.1346 & 3.33 \\
24 & 0.0465 & 0.0930 & 2.00 \\
\bottomrule
\end{tabular}

\begin{minipage}{0.85\textwidth}
\vspace{2pt}
{\scriptsize \emph{Notes:} Bootstrap standard error from the
exact-specification bootstrap; kernel standard error from the main
table. Source: \texttt{se\_ratio\_nb07.csv}.}
\end{minipage}
\end{table}

\paragraph{Block-length sensitivity.}
The 24-hour block length of the moving-block bootstrap is the maximal
horizon, the natural choice for the overlap structure. Nothing in our results
turns on it. Table~\ref{tab:blocksens} reports the 80\%-power minimum
detectable effect for the paired one-percent asymmetry, the object that
drives the equivalence bound, at block lengths of 12, 24, 36, and 48
hours. The bound moves by less than a tenth of its value across the range.
By construction the 24-hour column reproduces the corresponding cells of the
main exceedance table exactly, a built-in cross-check that the two run on
identical machinery.

\begin{table}[t]
\centering
\caption{Block-length sensitivity of the equivalence bound (MDE at 80\%
power, $\Delta$ at $\alpha = 1\%$).}
\label{tab:blocksens}
% GENERATED by scripts/paper/make_tables.py — DO NOT EDIT
% fragment only: wrap in \begin{table} + caption in the manuscript
\begin{tabular}{lcccc}
\toprule
$h$ & block $=12$h & block $=24$h & block $=36$h & block $=48$h \\
\midrule
0 & 0.00096 & 0.00094 & 0.00089 & 0.00091 \\
1 & 0.00101 & 0.00101 & 0.00096 & 0.00097 \\
3 & 0.00087 & 0.00087 & 0.00087 & 0.00088 \\
6 & 0.00113 & 0.00109 & 0.00097 & 0.00096 \\
12 & 0.00108 & 0.00114 & 0.00112 & 0.00114 \\
24 & 0.00102 & 0.00097 & 0.00099 & 0.00101 \\
\bottomrule
\end{tabular}

\begin{minipage}{0.85\textwidth}
\vspace{2pt}
{\scriptsize \emph{Notes:} Per-period paired bootstrap, $1{,}000$
replications. The 24-hour column reproduces the main exceedance results.
Source: \texttt{block\_sensitivity.csv}.}
\end{minipage}
\end{table}